% ============================================================================
% TWIN PRIMES AND SIEGEL ZEROS: A FULCRUM -- the PREPRINT build.
%
% main.tex is set in the Forum of Mathematics journal class (CUP-JNL-FMP.cls),
% which prints the journal's name, a DOI placeholder, the publisher's logos and
% a licence sentence on every page. That livery belongs on a submission to the
% journal and on nothing else, so the public preprint (the Zenodo record of
% 2026-10-02) is built from THIS file instead: the standard article class, a
% few definitions standing in for the journal's frontmatter macros, and then
% main.tex itself, unchanged, read with \input. Every sentence of the paper is
% main.tex's own; this file adds a title block and nothing else.
%
% Build: tectonic preprint.tex   (or pdflatex preprint.tex, twice)
% ============================================================================
\documentclass[10pt]{article}
% The journal class's own text block (CUP-JNL-FMP.cls: textwidth 144mm, textheight 612.5pt),
% so the Lean listings in Appendix A break where the journal build breaks them.
\usepackage[a4paper,textwidth=144mm,textheight=612.5pt]{geometry}
\usepackage{xcolor}
% Two of the journal class's typesetting settings, so long Lean names break as they do there:
% its paragraph tolerance, its 9pt \small and its 7pt \footnotesize, in which Appendix A's
% listings are set (CUP-JNL-FMP.cls: \tolerance=9999; \small 9/11; \footnotesize 7/8).
\tolerance=9999
\emergencystretch=3em
\makeatletter
\renewcommand\small{\@setfontsize\small\@ixpt{11}}
\renewcommand\footnotesize{\@setfontsize\footnotesize{7}{8}}
\makeatother
\definecolor{weblmcolor}{rgb}{0.0,0.25,0.55}

\makeatletter
% --- Stand-ins for the journal class's frontmatter macros (no output of their own).
\newcommand{\articletype}[1]{}
\newcommand{\jyear}[1]{}
\newcommand{\copytext}[1]{}
\newcommand{\orcid}[1]{\def\fp@orcid{#1}}
\newcommand{\authormark}[1]{}
\newcommand{\orgname}[1]{#1}
\newcommand{\fp@address}{}
\newcommand{\address}[2][]{\def\fp@address{#2}}
\let\fp@title\@empty
\renewcommand{\title}[2][]{\def\fp@title{#2}}
\renewcommand{\author}[2][]{\def\fp@author{#2}}
\newcommand{\fp@abstract}{}
\renewcommand{\abstract}[1]{\def\fp@abstract{#1}}
\newcommand{\fp@keywords}{}
\newcommand{\fp@msc}{}
\newcommand{\keywords}[2][]{%
  \ifx\relax#1\relax\def\fp@keywords{#2}\else\def\fp@msc{#2}\fi}
\newcommand{\codes}[2][]{#1~#2}
\newenvironment{Frontmatter}{}{%
  \begin{center}
    {\LARGE\bfseries \fp@title\par}\vspace{1.2em}
    {\large \fp@author\par}\vspace{0.3em}
    {\small \fp@address\par
      ORCID \href{https://orcid.org/\fp@orcid}{\fp@orcid}\par}\vspace{0.6em}
    {\small Preprint, 2 October 2026\par}
  \end{center}
  \vspace{0.8em}
  \begin{quote}\small
    \textbf{Abstract.} \fp@abstract\par\medskip
    \textbf{Keywords:} \fp@keywords.\par
    \textbf{MSC 2020:} \fp@msc.\par
  \end{quote}
  \vspace{0.8em}}
% --- main.tex opens with its own \documentclass line; this build has already
% chosen the class, so that line is read and discarded.
\renewcommand{\documentclass}[2][]{}
% main.tex loads hyperref; give the PDF its document-info title and author once it has.
\AtBeginDocument{\hypersetup{pdftitle={Twin Primes and Siegel Zeros: a Fulcrum},pdfauthor={Jason Hickey}}}
% The journal class prints its title block at \maketitle; here the Frontmatter
% environment has already printed it.
\renewcommand{\maketitle}{}
\makeatother

% ============================================================================
% TWIN PRIMES AND SIEGEL ZEROS: A FULCRUM
% v0.1 drafted 2026-07-19 (Dichotomy Day) per the noon council rulings
% (docs/exploration/council-brief-noon.md, D1-D6 + trust-surface + discovery).
% Target: publication approval -> arXiv (math.NT primary, cs.LO cross) -> Pi.
% Convention: every stated result carries its Lean name via \leanname{};
% Appendix A maps statement -> declaration -> axiom audit.
% Voice: problem-first; every strong adjective paid by a number; every
% conditional named at the statement. No local LaTeX toolchain on the build
% machine -- compile checks happen author-side/arXiv.
% ============================================================================
\documentclass[11pt]{amsart}
\usepackage[margin=1.1in]{geometry}
\usepackage{amsmath,amssymb,amsthm}
\usepackage{booktabs}
\usepackage[colorlinks=true,linkcolor=blue!60!black,citecolor=blue!60!black,urlcolor=blue!60!black]{hyperref}
\usepackage{microtype}

\newcommand{\leanname}[1]{\par\noindent{\small\texttt{#1}}}
\newcommand{\leaninline}[1]{{\small\texttt{#1}}}

\newtheorem{theorem}{Theorem}[section]
\newtheorem{corollary}[theorem]{Corollary}
\newtheorem{definition}[theorem]{Definition}
\newtheorem{proposition}[theorem]{Proposition}
\theoremstyle{remark}
\newtheorem{remark}[theorem]{Remark}

\title{Twin Primes and Siegel Zeros: a Fulcrum}
\author{Jason Hickey}
\thanks{Developed in collaboration with Claude (Anthropic's Fable~5 and
Opus models); see \S\ref{sec:method} and the repository's commit ledger,
in which every commit is co-authored. Written in a personal capacity.}
\date{July 19, 2026}

\begin{document}

\begin{abstract}
Every theorem in this paper is verified by the Lean~4 kernel over
\textsf{mathlib}, with axiom base exactly
$\{\texttt{propext}, \texttt{Classical.choice}, \texttt{Quot.sound}\}$;
nothing rests on informal argument. The mathematics---the designs, the
proofs, and their formalization---is joint work of the author and Claude
(Anthropic's Fable~5 and Opus models), carried out under the author's
direction and ratification; the commit ledger records the collaboration
line by line.

We introduce the \emph{fulcrum}: the weakest hypothesis yet formulated
under which Siegel zeros produce infinitely many twin primes---a single
fixed quality constant $C^\star$, one zero per unbounded window, with the
zero's reality \emph{derived} rather than assumed. The fulcrum was not
transcribed from the literature: it was extracted from a formal corpus by
demand-audited hypothesis minimization, a mode of discovery we argue is
impossible over prose proofs. Around it we prove, kernel-checked: a
Heath-Brown--type dichotomy (twin primes, or no Siegel zeros) with its
engine-conditional horn stated honestly; the first formalization of any
zero-free region for $\zeta$ of Littlewood strength (1922) and the first
power-saving region ($\theta = 3/4$) in any proof assistant; a
Deuring--Heilbronn repulsion contract; the exact-identity keystone of the
Heath-Brown engine, reduced to one named residual; the \emph{exchange-rate
wall}, a neutrality theorem over arbitrary real sign functions that prices
every transfer-based route to twins; and the \emph{parity gap} as a
mathematical object: a predicate $Z$ with a kernel-checked theorem that no
true, twin-sufficient, parity-invariant completion exists, together with a
census placing the entire corpus inside the parity-invariant cone. As a
byproduct, the logarithmic two-point Chowla conjecture is reduced to a
single named hypothesis. The corpus is $317{,}039$ lines of Lean across
$732$ files; the method that produced it---the \emph{Salt method}---kept a
public ledger of every design error caught (252 numbered entries) against
zero wrong proofs, and is described as a contribution in its own right.
\end{abstract}

\maketitle

% ============================================================================
\section{Introduction}\label{sec:intro}
% ============================================================================

The twin prime conjecture is not merely unproven; the routes toward it are
walled by a structural obstruction---the parity barrier---that has never
itself been a theorem, and guarded by a ghost---the Siegel zero---that has
never been an object one could bargain with precisely. Progress of the
conditional kind (Heath-Brown's 1983 theorem that infinitely many Siegel
zeros force infinitely many twin primes \cite{HB1983}) trades in
hypotheses whose exact strength is scattered through pages of estimates:
what, \emph{precisely}, must one believe about exceptional zeros to
believe in twin primes? Prose proofs cannot answer such questions
exhaustively, because a prose proof hides its demand.

A formal proof \emph{is} its demand. This paper's central object, the
\textbf{fulcrum}, was found by exploiting exactly that: we formalized a
Heath-Brown--type engine, enumerated---mechanically, exhaustively---every
site at which the Siegel-zero hypothesis is consumed and at what strength,
and then weakened the hypothesis axis by axis to the audited demand, with
the Lean kernel refereeing each weakening instantly. The fixed point of
that process is the weakest twin-producing Siegel hypothesis yet
formulated (Definition~\ref{def:fulcrum}); the process itself,
\emph{demand-audited hypothesis minimization}, is to our knowledge a new
mode of mathematical discovery, and it is the sense in which this paper's
formalization is an instrument rather than a notary.

The results divide into four movements. First, the \emph{banners}
(\S\ref{sec:banners}): unconditional analytic theorems formalized for the
first time, including a Littlewood-strength zero-free region for $\zeta$
and the first power-saving region in any proof assistant. Second, the
\emph{fulcrum and the dichotomy} (\S\S\ref{sec:fulcrum}--\ref{sec:keystone}):
the minimized hypothesis, the kernel-checked dichotomy it induces, and the
exact-identity keystone of the engine, brought to within one named lemma
of closing. Third, the \emph{negative space}
(\S\S\ref{sec:walls}--\ref{sec:parity}): the exchange-rate wall and the
parity gap, which together map---as theorems, not folklore---why the
obvious routes die. Fourth, the \emph{method} (\S\ref{sec:method}): the
Salt method, an adversarial division of mathematical labor among human
direction, frontier-model design, fleet-model execution and refutation,
and kernel verification, whose public error ledger (252 numbered catches,
zero wrong proofs) is offered as evidence that the standard of the first
sentence of the abstract is achievable at research scale and research
speed. All of it was produced by the Salt method (\S\ref{sec:method}).

Throughout, conditionality is a first-class citizen: every conditional
statement names its exact residual hypothesis, and
Table~\ref{tab:conditional} lists all of them in one place. The paper
makes no unconditional claim about twin primes.

% ============================================================================
\section{The corpus and the referee}\label{sec:trust}
% ============================================================================

The corpus is $317{,}039$ lines of Lean~4 over \textsf{mathlib}
\cite{mathlib}, in $732$ files, built by $999$ commits; a full build
kernel-checks roughly $8{,}800$ compilation jobs. The skeptic's question
---\emph{how do we really know?}---splits into what is mechanically
excluded and what is irreducibly a human question, and we answer it in
that order.

\subsection{Mechanically excluded}
A \texttt{sorry} anywhere in a dependency tree---ours or
\textsf{mathlib}'s---introduces the axiom \texttt{sorryAx}, which
propagates transitively and surfaces in the axiom closure of every
downstream theorem. Our \texttt{\#audit\_axioms} blocks print that closure
for every headline statement: it is exactly
$\{\texttt{propext}, \texttt{Classical.choice}, \texttt{Quot.sound}\}$,
the three classical axioms of \textsf{mathlib} practice. A hidden gap
cannot hide. Tactic bugs cannot certify falsehoods: tactics emit proof
terms that the kernel re-checks independently, so the trust surface is the
kernel alone, a deliberately small core; we additionally forbid
\texttt{native\_decide} (which would trust the compiler, visibly, as a
fourth axiom). The kernel itself is the one residual mechanical trust,
and it is checkable: before submission the build is run through an
independent checker (\texttt{lean4checker}); two unrelated implementations
agreeing shrinks the kernel-bug hypothesis below the noise floor of any
classical refereeing process.

\subsection{The specification gap}
What no formal system can check is whether our \emph{statements} mean what
our prose claims. This gap is irreducible---it is the map--territory
boundary---so we engineer against it. Definitions bottom out in
\textsf{mathlib}'s communal objects (\texttt{Nat.Prime},
\texttt{DirichletCharacter}, \texttt{LFunction}), vetted by thousands of
independent uses. Vacuity is an attacked property, not an assumed one:
adversarial review passes carry mandatory degenerate-instance checks, and
where a definition provably trivializes outside its intended window, the
degeneracy is \emph{stated in Lean} beside it (see the grade guard of
\S\ref{sec:parity}). Non-vacuity is proved by witness: the fulcrum's
zero-reality is derived from our own zero-free region
(Theorem~\ref{thm:fulcrumreal}); the parity census instantiates its cone
with nine corpus theorems; classical corollaries are re-derived from our
statements, which a misstated premise would break. Finally,
Appendix~\ref{app:index} renders every headline statement in Lean and in
prose side by side, reducing the entire correctness question of the corpus
to a finite, public, refereeable translation audit. ``Absolutely'' is not
on offer anywhere in mathematics; what is on offer here is a doubt surface
consisting of one small kernel, three axioms, and a one-page table.

% ============================================================================
\section{The four banners}\label{sec:banners}
% ============================================================================

Four unconditional theorems anchor the corpus; none had been formalized
before, and two had never been formalized at any strength in any proof
assistant.

\begin{theorem}[Littlewood-strength zero-free region, 1922]\label{thm:litt}
There exist effective constants such that $\zeta(s) \neq 0$ in a region of
Littlewood type
$\sigma \ge 1 - c\,\frac{\log\log |t|}{\log |t|}$ for $|t|$ large.
\leanname{Salt.Vk.zeta\_zero\_free\_region\_littlewood}
\end{theorem}

\begin{theorem}[The first power-saving region]\label{thm:pow}
There exist effective constants such that $\zeta(s) \neq 0$ for
$\sigma \ge 1 - c\,(\log|t|)^{-3/4}(\log\log|t|)^{-O(1)}$, $|t|$ large:
a Vinogradov--Korobov--type region at exponent $\theta = 3/4$.
\leanname{Salt.Vk.zeta\_zero\_free\_region\_pow}
\end{theorem}

To our knowledge Theorem~\ref{thm:pow} is the first zero-free region with
a power saving in the exponent verified in any proof assistant, and
Theorem~\ref{thm:litt} the first formalization of the 1922 Littlewood
region; the strongest previously formalized regions were of classical
de la Vall\'ee Poussin shape. The proofs formalize the Vinogradov mean
value theorem (\leaninline{Salt.Vmvt.vmvt}) and route it through an
exponential-sum-to-growth bridge (\leaninline{Salt.Vk.vk\_block\_core}),
with an explicit strip family at absolute constant $C = 4096$.

\begin{theorem}[Deuring--Heilbronn repulsion, contract form]\label{thm:rep}
With explicit witnesses $b = 680$, $k = 14$, $\sigma_0 = 16/17$ and an
explicit constant $c = c_1^8$: an exceptional real zero close to $s=1$
repels every other zero of every Dirichlet $L$-function of comparable
modulus, in the ordered two-zero form.
\leanname{Salt.SW.dh\_repulsion\_ordered}
\end{theorem}

\begin{theorem}[Vinogradov mean value theorem]\label{thm:vmvt}
The mean value estimate for Weyl sums in the range needed by
Theorem~\ref{thm:pow}, with explicit constants.
\leanname{Salt.Vmvt.vmvt}
\end{theorem}

Around these sit an unconditional Siegel--Walfisz theorem, a
Bombieri--Vinogradov theorem, a machine-checked Chen's theorem with the
unconditional corollary $\Omega(p(p+2)) \le 3$ for infinitely many primes
$p$ (\leaninline{Salt.Fulcrum.chen\_omega\_prod\_le\_three}), and the
unconditional $\Omega(n(n+2)) \le 20$ infinitely often
(\leaninline{Salt.BrunLower.twin\_almost\_prime}); these are inherited
infrastructure from the same program and are indexed in
Appendix~\ref{app:index}.

% ============================================================================
\section{The fulcrum}\label{sec:fulcrum}
% ============================================================================

\subsection{The discovery}
The engine of \cite{HB1983} consumes a Siegel-zero hypothesis at many
sites: as a lower bound forcing $L$-values small, as a localization of one
zero, as a quality bound entering a sieve level, as a reality assumption.
In prose, the totality of this demand is not enumerable; in the formal
corpus it is a finite list of typed consumption sites, each with an exact
strength. The Fulcrum Hunt (three adversarial passes over the corpus,
recorded in the repository's exploration ledger) audited every site,
then weakened the hypothesis along each axis until the audited demand was
met exactly---each candidate weakening compiled or refuted by the kernel
within minutes. Three collapses survived: the quality constant collapsed
from a universally quantified family to a \emph{single fixed} $C^\star$;
the zero's \emph{reality} moved from hypothesis to conclusion (derived
from our own zero-free region); and the \emph{quantity} collapsed to one
zero per unbounded window---no asymptotic family, no density, no
infinitude of characters beyond what one zero per window forces.

\begin{definition}[The fulcrum]\label{def:fulcrum}
For $C > 0$, \emph{FulcrumQualityMin} $C$ asserts: for every $Q$ there
exist a modulus $q > Q$, a primitive quadratic character $\chi \bmod q$,
and a zero $\rho$ of $L(\cdot,\chi)$ with
$\|1-\rho\| \cdot C \log q \le 1$.
The fulcrum hypothesis is FulcrumQualityMin $C^\star$ for the single
explicit constant $C^\star = \max(C^{(1)}, 2/c_0)$, where
$c_0 = 1/126848$ is the zero-free-region constant of the corpus,
kernel-certified at that numeral
(\leaninline{zero\_free\_region\_all\_numeral}); the $2/c_0$ arm of
$C^\star$ is exactly $253696$, and the fulcrum's reality theorem is
instantiated at that bare threshold
(\leaninline{fulcrum\_zero\_real\_numeral}).
\leanname{Salt.Fulcrum.FulcrumQualityMin}
\end{definition}

\begin{theorem}[Reality is derived]\label{thm:fulcrumreal}
Any zero witnessing the fulcrum at quality $C^\star$ is real, and lies in
the exceptional position of Siegel-zero theory.
\leanname{Salt.Fulcrum.fulcrum\_zero\_real}
\end{theorem}

\begin{remark}
One Siegel zero is one bounded window: finding a single exceptional zero
would refute GRH but would \emph{not} satisfy the fulcrum, which demands a
witness beyond every $Q$. The fulcrum carries no asymptotic content
per window; this is the weakest shape for which the engine of
\S\ref{sec:keystone} can run, and the demand audit certifies that
\emph{this} engine leaves nothing weaker on the table.
\end{remark}

% ============================================================================
\section{The dichotomy}\label{sec:dichotomy}
% ============================================================================

\begin{theorem}[The dichotomy, engine-conditional horn stated
honestly]\label{thm:dich}
Unconditionally: if the fulcrum fails, there are no Siegel zeros in the
precise sense of the corpus's \leaninline{NoSiegelZeros}
(\leaninline{Salt.Fulcrum.not\_fulcrum\_implies\_noSiegelZeros}).
Consequently, for any $C > 0$, given the engine implication
$h_{\mathrm{Engine}} : \text{FulcrumQualityMin } C \to
\text{TwinPrimeConjecture}$, we obtain
\[
\text{TwinPrimeConjecture} \ \lor\ \text{NoSiegelZeros}.
\]
\leanname{Salt.Fulcrum.fulcrum\_dichotomy}
\end{theorem}

The second horn is the Heath-Brown dichotomy, machine-checked and
strengthened: strengthened because the hypothesis feeding its twin horn is
the minimized fulcrum rather than the 1983 hypothesis, and because the
$\lnot F$ horn is unconditional and effective. The engine implication
$h_{\mathrm{Engine}}$ is the one conditional input; \S\ref{sec:keystone}
reports its keystone proven and its remaining distance reduced to one
named lemma. We do not claim the unconditional dichotomy; we claim the
honest one, and we exhibit exactly what remains.

% ============================================================================
\section{The keystone}\label{sec:keystone}
% ============================================================================

The engine's hardest block is the second-moment comparison
$S_2 - S_1$ over the sifted twin window---the site of \cite{HB1983}'s
Lemma~2. Formalization exposed that the classical route's
$\tau$-crude majorant is \emph{provably} insufficient at the needed
generality (its concentration pattern is unexcludable by pattern-blind
tools; ledger entry \#251), so the corpus proves the required conclusion
directly on an exact identity: the overshoot
$(\tilde\Lambda-\Lambda)(n)\tilde\Lambda(n+2) +
\Lambda(n)(\tilde\Lambda-\Lambda)(n+2)$ is summed exactly, its support is
classified, and eleven family estimates (each with explicit absolute
constants; nine landed at first attempt) price every class.

\begin{theorem}[The master estimate]\label{thm:master}
With $W$ the sifted window at parameters
$(z,x)$ in the engine's regime, and with one named residual hypothesis
$h_{\mathrm{count}}$ (below):
\[
S_2(\chi, W) - S_1(W) \;\le\; 2^{31}\Bigl(\frac{x}{z_0}
+ \frac{x}{\log x}\,e^{5 z_0}\,\mathrm{PS}_\chi
+ e^{2 z_0}\bigl(x z^{-1/8} + x^{9/10}\bigr) L^3\Bigr),
\]
where $z_0 = \log(2x+2)/\log z$, $L = \log(2x+2)$, and
$\mathrm{PS}_\chi$ is the pretense sum of $\chi$.
\leanname{Salt.HB.hb\_l2c\_master\_final}
\end{theorem}

The residual $h_{\mathrm{count}}$ is a $\chi$-weighted sieve count (the
roles-swapped family of the mirror side, where the character-positive
prime is a cofactor rather than a modulus); at the engine's own witness
scale one boundary family requires a second count of the same shape, and
both are stated exactly in the corpus and registered together as the
single node \textsc{chi-sieve}---the engine's entire remaining
sieve-theoretic content. The cover of the eleven families is itself a
kernel-checked theorem---no case can have been forgotten---and the three
cover gaps found during execution were all found \emph{by the machine
layers below the kernel} and repaired at the statement level before any
proof was attempted on a wrong statement (ledger \#245, \#246, \#252).
The theorem as stated holds in the $z = x^{o(1)}$ regime, where its full
hypothesis packet is discharged by explicit witnesses
(\leaninline{Salt.HB.glue\_master}); the regime bookkeeping itself
produced one of the campaign's sharpest catches, an amendment refuted by
its own executor at the asymptotic corner (ledger \#253).

% ============================================================================
\section{The exchange-rate wall}\label{sec:walls}
% ============================================================================

The Fulcrum Hunt's second product is negative space made precise: four
recorded refutations (the \emph{wall series}, banked in the repository's
exploration ledger) of the natural routes by which Siegel-zero
information might be converted into twin-prime information. Three are
ledger-grade records; the fourth is a theorem series, formalized here.

The transfer at the heart of every such route exchanges the twin sum
$S_1$ for a twisted sum $S_2(f)$ against a multiplicative sign function
$f$. The wall states: this exchange is \emph{never free, and its price is
exactly a pretense sum}---and it says so for arbitrary real sign
functions, a class strictly wider than Dirichlet characters.

\begin{definition}[Sign functions]\label{def:sign}
$f : \mathbb{N} \to \mathbb{R}$ is a \emph{sign function} if it is
totally multiplicative with $f(1)=1$, $|f| \le 1$, and $f(p) = \pm 1$ on
primes.
\leanname{Salt.HB.IsSignFunction}
\end{definition}

\begin{theorem}[Neutrality and the rate]\label{thm:wall}
For every sign function $f$: $S_1 \le S_2(f)$ on the generalized window
(unconditionally), and the difference obeys the budget
$S_2(f) - S_1 \le \mathrm{EngineBound}(C_{\mathrm{main}}, z, x,
\mathrm{PS}_f)$ whenever the overshoot budget holds---the three-row bound
of Theorem~\ref{thm:master} with the character's pretense sum replaced by
$f$'s.
\leanname{Salt.HB.S1\_le\_S2Gen, Salt.HB.neutrality\_rate}
\end{theorem}

\begin{theorem}[The neutrality point]\label{thm:liouville}
The Liouville function is the exact fixed point of the exchange:
$\tilde\Lambda_\lambda = \Lambda$ identically, $S_2(\lambda) = S_1$, the
overshoot vanishes pointwise, and $\mathrm{PS}_\lambda = 0$.
\leanname{Salt.HB.LamTildeGen\_lamR\_eq\_vonMangoldt,
S2Gen\_lamR\_eq\_S1, overshootExactGen\_lamR, PretenseSumGen\_lamR\_eq\_zero}
\end{theorem}

Read together: any route that would extract twins from a transfer must pay
in pretense, the currency $\mathrm{PS}_f$; the one sign function with
$\mathrm{PS} = 0$ is the one at which the transfer returns exactly the
twin sum it was fed. The wall is falsifiable---exhibit a sign function
with bounded pretense sum and an $S_2$-evaluation that does not route
through $S_1$---and it is stated over the full sign-function cone
precisely so that its scope is a theorem rather than an impression. The
re-typing of the entire transfer chain from Dirichlet characters to sign
functions (59 declarations) required deleting every coprimality thread;
that the deletion is sound is itself part of the formal content: the
character was never load-bearing.

% ============================================================================
\section{The parity gap, as an object}\label{sec:parity}
% ============================================================================

The parity barrier has lived in the literature as folklore with
heuristics. The corpus makes it an object with a theorem.

\begin{definition}[Completions, parity invariance, and $Z$]\label{def:z}
Over weightings $a$ of the integers, a \emph{completion} at grade
$(\theta, A_0)$ is a predicate on weightings capturing type-I control at
level $x^\theta$ with logarithmic savings $A_0$; \emph{ParityInv} is the
cone of completion-predicates invariant under the parity involution;
\emph{TwinSufficient} demands that a predicate force unbounded twin mass;
and $Z(\theta, A_0)$ asserts the existence of a true, twin-sufficient
completion predicate. The file defining these
(\leaninline{Salt/Parity/Z.lean}) imports nothing from the Siegel-zero or
sieve tracks: its oracle-cleanliness is checkable by import list alone.
\leanname{Salt.Parity.Z, Salt.Parity.ParityInv}
\end{definition}

\begin{theorem}[The gap theorem]\label{thm:gap}
No completion predicate is simultaneously true, twin-sufficient, and
parity-invariant: the twin-free weighting lies in every parity-invariant
completion cone at the certified grades, and its twin mass is identically
zero.
\leanname{Salt.Parity.sufficient\_true\_not\_parityInv}
\end{theorem}

\begin{theorem}[$Z$ is exactly the demand]\label{thm:ziff}
Over the certified grade window ($0 < \theta < 1/2$, $A_0$ in the stated
range): $Z(\theta,A_0) \iff$ the twin prime conjecture.
\leanname{Salt.Parity.Z\_implies\_TPC, Salt.Parity.TPC\_implies\_Z}
\end{theorem}

\begin{proposition}[The census]\label{prop:census}
Nine headline theorems of the corpus---including the $k=2$ Maynard bound
$M_2 \le 2\log 2$, the Chen-style headline, the Brun-grade twin bound, and
the sieve keystones---lie inside the parity-invariant cone.
\leanname{Salt/Parity/Instances.lean (nine declarations)}
\end{proposition}

The three results triangulate the frontier exactly: everything the corpus
proves lives inside the cone (Proposition~\ref{prop:census}); the twin
prime conjecture is equivalent to a demand $Z$
(Theorem~\ref{thm:ziff}); and that demand cannot be met inside the cone
(Theorem~\ref{thm:gap}). The barrier is now a boundary with coordinates.
Where $Z$ degenerates outside its certified window, the degeneracy is
itself a Lean lemma beside the definition (the grade guard,
\leaninline{Z\_trivial\_of\_not\_completion})---vacuity disclosed, not
discovered.

% ============================================================================
\section{The door-only spine}\label{sec:spine}
% ============================================================================

Independently of the Siegel-zero tracks, the corpus carries a full
formalization of the entropy-decrement route to the logarithmic two-point
Chowla conjecture. On the day of writing, its residual collapsed: the
prior interface (three coupled analytic witnesses) was \emph{proven
unsatisfiable} at its recorded constants---a freeze-panel finding,
confirmed four independent ways, ledger \#249---and the corrected
budget was proven through, leaving:

\begin{theorem}[Log-Chowla, door-only]\label{thm:spine}
There exist $\delta_0 > 0$ and a regime $R$ such that a single hypothesis
---the Matom\"aki--Radziwi\l\l--Tao uniformity door
$\mathrm{MRTUniformity}\,R\,\delta$ at any $0 < \delta \le \delta_0$---
implies the logarithmic two-point Chowla statement at $R$. A
$\forall$-quantified head version places the regime floor above any
prescribed threshold.
\leanname{Salt.Entropy.Chowla.log\_chowla\_two\_door\_only,
log\_chowla\_two\_budget\_head}
\end{theorem}

The door is the exact target of the corpus's staged
Matom\"aki--Radziwi\l\l\ program (whose first stones, including the
closure of a long-flagged Euler-oscillation bridge, landed the same day);
the conjecture's entire residual is thus one named seam between two
formalized campaigns.

% ============================================================================
\section{The Salt method}\label{sec:method}
% ============================================================================

\begin{quote}
\emph{The Salt method is an adversarial division of mathematical labor:
human direction and ratification; frontier-model design under statement
freezes; fleet-model execution and independent refutation, tiered by
difficulty class; and kernel verification as the sole referee. Its
discipline---freeze before proving, refute before trusting, give up
loudly, bank everything---is designed to make the system catch errors
faster than it makes them.}\footnote{After water, salt.}
\end{quote}

Every lemma is classified before any proof attempt (one-step; standard;
real design; open-ended) and routed to the cheapest layer that can carry
it; statements are frozen by design panels whose members include
adversarial refuters with a standing default of \emph{refuted}; executors
are forbidden to alter a frozen statement---their duty on discovering one
unprovable is to stop and report, loudly; and the house that rules on such
reports is itself subject to the same audit (twice in the campaign, an
executor refuted the supervising layer's own amendment, and the record
says so: ledger \#248, \#250).

The evidence that this discipline works is the ledger: 252 numbered
catches over the campaign against \emph{zero} wrong proofs built on any of
them. The day this paper was drafted supplies the case study: five
statement-layer faults surfaced in one morning---three cover-taxonomy
gaps in a frozen design, one refutation of a supervising amendment, one
unsatisfiable recorded interface---and every one was caught by the layers
below the kernel, ruled on at the statement level, and repaired before a
single wrong proof was attempted. A companion paper (\emph{The Salt
Method}) will give the full treatment; the repository's ledgers are the
raw artifact both papers cite.

We record one measurement with a claim attached: at no point in the
campaign did the kernel discover an error. It confirmed; the adversarial
layers discovered. A verification standard is achievable at research
speed precisely because correctness is enforced \emph{above} the kernel
and only certified at it.

% ============================================================================
\section{The program forward}\label{sec:forward}
% ============================================================================

Four named campaigns continue from the corpus's current edge, each with
its target stated formally in advance: \textsc{hb-engine} (the remaining
distance from Theorem~\ref{thm:master} to $h_{\mathrm{Engine}}$);
\textsc{chi-sieve} (the one residual of the keystone); the
Matom\"aki--Radziwi\l\l\ program toward the door of
Theorem~\ref{thm:spine}; and the explicit-constants floor
(Siegel--Walfisz through Chen with every constant explicit). The
conditional ledger below is the paper's honest boundary.

\begin{table}[h]
\centering\small
\begin{tabular}{lll}
\toprule
Conditional statement & Residual hypothesis & Campaign \\
\midrule
Thm~\ref{thm:dich}, twin horn & $h_{\mathrm{Engine}}$ & \textsc{hb-engine} \\
Thm~\ref{thm:master} & $h_{\mathrm{count}}$ & \textsc{chi-sieve} \\
Thm~\ref{thm:spine} & the MRT door & S8 / MR program \\
Euler-oscillation bridge & $\mathrm{PrimeTailShiftBounded}$ & R0.2 \\
\bottomrule
\end{tabular}
\caption{Every conditional in this paper, with its exact named residual.
Nothing else in the paper is conditional.}
\label{tab:conditional}
\end{table}

% ============================================================================
\section*{Acknowledgments}
% ============================================================================
This work was carried out in a personal capacity. The formalization
builds on \textsf{mathlib} and its community. The collaboration with
Claude (Anthropic's Fable~5 and Opus models) is described in
\S\ref{sec:method} and recorded commit-by-commit in the repository; the
author directed and ratified every design decision and bears sole
responsibility for the claims.

% ============================================================================
\appendix
\section{The formal-statement index}\label{app:index}
% ============================================================================
% At release: generated from the All.lean audit blocks; every row's axiom
% column verified by #audit_axioms. Prose <-> Lean side-by-side renderings
% for the headline statements follow the table. [v0.1: headline rows only]

\begin{table}[h]
\centering\small
\begin{tabular}{lll}
\toprule
Statement & Declaration & Axioms \\
\midrule
Thm~\ref{thm:litt} & \leaninline{zeta\_zero\_free\_region\_littlewood} & the three \\
Thm~\ref{thm:pow} & \leaninline{zeta\_zero\_free\_region\_pow} & the three \\
Thm~\ref{thm:rep} & \leaninline{dh\_repulsion\_ordered} & the three \\
Def~\ref{def:fulcrum} & \leaninline{FulcrumQualityMin} & --- \\
Thm~\ref{thm:fulcrumreal} & \leaninline{fulcrum\_zero\_real} & the three \\
Thm~\ref{thm:dich} & \leaninline{fulcrum\_dichotomy} & the three \\
Thm~\ref{thm:master} & \leaninline{hb\_l2c\_master\_final} & the three \\
Thm~\ref{thm:wall} & \leaninline{neutrality\_rate} & the three \\
Thm~\ref{thm:liouville} & \leaninline{S2Gen\_lamR\_eq\_S1} & the three \\
Thm~\ref{thm:gap} & \leaninline{sufficient\_true\_not\_parityInv} & the three \\
Thm~\ref{thm:ziff} & \leaninline{Z\_implies\_TPC / TPC\_implies\_Z} & the three \\
Thm~\ref{thm:spine} & \leaninline{log\_chowla\_two\_door\_only} & the three \\
\bottomrule
\end{tabular}
\caption{``The three'' $=$ \leaninline{propext, Classical.choice,
Quot.sound}. Full index generated at release.}
\end{table}

\begin{thebibliography}{99}
\bibitem{HB1983} D.~R. Heath-Brown, \emph{Prime twins and Siegel zeros},
Proc.\ London Math.\ Soc.\ (3) \textbf{47} (1983), 193--224.
\bibitem{Chen1973} J.~R. Chen, \emph{On the representation of a larger
even integer as the sum of a prime and the product of at most two primes},
Sci.\ Sinica \textbf{16} (1973), 157--176.
\bibitem{Littlewood1922} J.~E. Littlewood, \emph{Researches in the theory
of the Riemann $\zeta$-function}, Proc.\ London Math.\ Soc.\ (2)
\textbf{20} (1922).
\bibitem{TaoChowla} T.~Tao, \emph{The logarithmically averaged Chowla and
Elliott conjectures for two-point correlations}, Forum Math.\ Pi
\textbf{4} (2016).
\bibitem{MR2016} K.~Matom\"aki and M.~Radziwi\l\l, \emph{Multiplicative
functions in short intervals}, Ann.\ of Math.\ \textbf{183} (2016).
\bibitem{MRT} K.~Matom\"aki, M.~Radziwi\l\l, and T.~Tao, \emph{An
averaged form of Chowla's conjecture}, Algebra Number Theory \textbf{9}
(2015).
\bibitem{Maynard2015} J.~Maynard, \emph{Small gaps between primes},
Ann.\ of Math.\ \textbf{181} (2015), 383--413.
\bibitem{HalesKepler} T.~Hales et al., \emph{A formal proof of the Kepler
conjecture}, Forum of Mathematics, Pi \textbf{5} (2017), e2.
\bibitem{mathlib} The mathlib community, \emph{The Lean mathematical
library}, CPP 2020.
\bibitem{companion} J.~Hickey, \emph{The Salt method}, in preparation.
\end{thebibliography}

\end{document}

